% Exposé VII. Worked examples: fourteen methods through the three invariants
% Sources: theory/framework.md (Exposé VII), site/sections/10-atlas.html (#atlas-examples: the fourteen cards and their
% credit note), public statements in theory/propositions.json, and the rows of framework.md VIII.B for these methods.
% No proofs here: every claim is a reading of a numbered result proved in Appendix C.
\section{Worked examples}\label{sec:examples}

% local layout for the fourteen cards (names carry the VII tag to stay unique in main.tex)
\newcommand{\exVIIlabel}[1]{{\sffamily\footnotesize\bfseries\color{inkgray}\textls[80]{\MakeUppercase{#1}}}}
\newcommand{\exVIIsrc}[2]{{\small\color{inkgray}#1\enspace\textperiodcentered\enspace{\sffamily\footnotesize #2}}}
\newcommand{\exVIIrefs}[1]{\par\smallskip\noindent{\footnotesize\color{inkgray}{\sffamily\bfseries Results.}\ #1}}
\newenvironment{exVIIlines}
  {\begin{description}[leftmargin=0pt,labelindent=0pt,itemsep=3pt,topsep=2pt,font=\exVIIlabel]}
  {\end{description}}

This exposé reads fourteen methods through the three invariants and the merge layer. For each method it says what the
method reaches, how training moves it, what a moved base point does and how it merges, and it names the numbered
result behind every claim.

A method is a pointed map $\rho:(Q,q_0)\to(\Theta,\theta_0)$. Its image germ says what it can express, its fibres
together with the optimizer's metric say how it moves, base change says what a moved $\theta_0$ does, and extensions
cover adapters and prompts. The fourteen cards below apply that toolkit; Exposé~\ref{sec:atlas} turns the answers into
coordinates and arrows.

Each card gives the object $(Q,q_0,\rho)$ for one linear layer $y=Wx$ with $W\in\R^{m\times n}$, and then reads it
through the invariants. The \exVIIlabel{Image} line is the image germ at $\theta_0$ (Exposé~\ref{sec:image}); the
\exVIIlabel{Fibre} line is the fibre, the gauge and the induced metric (Exposés~\ref{sec:fibres} and~\ref{sec:lens});
the \exVIIlabel{Base} line is base change, meaning pointings, quantised bases and merge-and-restart
(Exposé~\ref{sec:base}); and the \exVIIlabel{Merge} line is the extension and the merge type
(Exposé~\ref{sec:merge}). The prefix card adds an \exVIIlabel{Arrows} line. Lines that an invariant leaves empty or
trivial are omitted. The merge types are those of Exposé~\ref{sec:atlas}, where M1 merges into its own box, M1$\to$ into a
neighbouring box under side conditions, Mq leaves the quantised type when merged exactly, and M$\infty$ is not
box-locally mergeable. The header of each card names the source paper and the modification kind used in the
catalogue, and its last line lists the results its claims rest on. Their complete proofs are in
Appendix~\ref{app:proofs}; claims marked [Num] are numerical checks (Appendix~\ref{app:repro}).

\paragraph{Credits}
Balancedness is due to \citet{du2018algorithmic} and \citet{arora2018optimization}, conserved charges to
\citet{zhao2022symmetries}, and the balanced slice to \citet{xkempf1979length}. The crossover is the closed form of
\citet{xtarmoun2021understanding} and \citet{xmin2021explicit}. The comparison of zero-$B$ and zero-$A$ starts is due to
\citet{hayou2024impact}, and scaled GD to \citet{xtong2021accelerating} and \citet{zhang2024riemannian}. GaLore as
one-sided LoRA with carried state is due to \citet{xtorrobahennigen2025duality}. Prefixes as gated adapters are due to
\citet{he2021unified}, and the rigid content ratios and the prompt-to-prefix map to \citet{petrov2023prompting}. The
merge of (IA)$^3$ into $W_{\rm down}$ is in \citet{liu2022few}, and scale migration through SwiGLU is used by AWQ
\citep{xlin2024awq}. LoKr's transport is the rearrangement of \citet{loan1993approximation}. The $2r$ bound underlies
the PiSSA-to-LoRA conversion of \citet{meng2024pissa}.

% ------------------------------------------------------------------------------------------------ LoRA
\begin{example*}{LoRA}\phantomsection\label{ex:lora}%
\exVIIsrc{\citet{hu2021lora}}{additive}
\[
  \rho(B,A)=W_0+sBA,\qquad q_0=(0,A_0),\qquad B\in\R^{m\times r},\ A\in\R^{r\times n}.
\]
\begin{exVIIlines}
\item[Image] The cone $W_0+\Mr$, of dimension $r(m+n-r)$ for $r<\min(m,n)$. With $\rk A_0=r$ the first-order image is
  $S_1=\{XA_0\}$, of dimension $mr$, so $\theta_0$ is an apex with deficit $r(n-r)$ (Figure~\ref{fig:apex}a).
\item[Fibre] Gauge $\GL_r$, acting by $(B,A)\mapsto(Bg^{-1},gA)$. Under gradient flow with block learning rates
  $\eta_B,\eta_A$ and no dropout, the induced cometric is $K(G)=s^2(\eta_B\,GA\T A+\eta_A\,BB\T G)$. The Lie algebra splits as $\gl_r=\so_r\oplus\Sym_r$. The
  isometric half $\so_r$ leaves $K$ unchanged, and the $\Sym_r$ half carries the charge
  $\Phi=B\T B/\eta_B-AA\T/\eta_A$, conserved under gradient flow and equal to $-A_0A_0\T/\eta_A$ from the zero-$B$
  start. For isotropic $\Phi=-\kappa I$ the flow is solved in closed form
  \citep{xtarmoun2021understanding,xmin2021explicit}. With $\mu=s\sqrt{\eta_A\eta_B}$ and the crossover scale
  $\sigma^*=\mu\lvert\kappa\rvert/2$, zero-$B$ LoRA is one-sided well below $\sigma^*$, where effectively only $B$ moves
  as in LoRA-FA, and well above it follows the balanced dynamics of \citet{arora2018optimization}
  (Figure~\ref{fig:noether}). On the full-rank locus, with channel-uniform hyperparameters, SGD keeps $\Ort(r)$ of the
  gauge, Adam and AdamW keep the signed permutations $B_r$, and undamped scaled GD keeps all of $\GL_r$
  (Figure~\ref{fig:gauge}). Under Adam with $\varepsilon=0$, no weight decay, no gradient clipping and a learning-rate
  schedule shared by both blocks, LoRA+ with learning-rate ratio $\lambda$ follows exactly the trajectory of LoRA with
  $\alpha$ multiplied by $\lambda$ \citep{schulman2025lora} (Figure~\ref{fig:orbit}).
\item[Base] $K$ merge-and-restart cycles reach exactly $W_0+\MM_{\le\min(Kr,m,n)}$, so all of $\Theta$ at
  $K^*=\lceil\min(m,n)/r\rceil$. Exact full-rank pointings form three strata (zero $B$, zero $A$, and a split that
  subtracts $B_0A_0$ from the frozen weight). Their isomorphism classes are fixed by $\operatorname{row}A_0$,
  $\operatorname{col}B_0$ and $B_0A_0$ respectively, so as a set they are $\Gr(r,n)\sqcup\Gr(r,m)\sqcup\MM_r$.
\item[Merge] M1, since the trained update $sBA$ adds to $W_0$.
\end{exVIIlines}
\exVIIrefs{\thmref{Theorem}{II.2} \textperiodcentered{} \thmref{Proposition}{II.3} \textperiodcentered{}
  \thmref{Observation}{III.2} \textperiodcentered{} \thmref{Theorem}{III.5} \textperiodcentered{}
  \thmref{Corollary}{III.6} \textperiodcentered{} \thmref{Theorem}{III.9} \textperiodcentered{}
  \thmref{Corollary}{III.11} \textperiodcentered{} \thmref{Theorem}{III.12} \textperiodcentered{}
  \thmref{Theorem}{V.3}}
\end{example*}

% ------------------------------------------------------------------------------------------------ DoRA
\begin{example*}{DoRA}\phantomsection\label{ex:dora}%
\exVIIsrc{\citet{liu2024dora}}{hybrid}
\[
  \rho(\mu,B,A)=\diag(\mu)\,N(W_0+sBA),\qquad q_0=(\mu_0,0,A_0),\qquad \mu_{0,i}=\lVert w_{0,i}\rVert.
\]
Here $N$ normalises each row to unit length and $w_{0,i}$ is the $i$-th row of $W_0$.
\begin{exVIIlines}
\item[Image] For $W_0$ without zero rows and $r\le\min(m,n)$ the image is
  $\mathrm{Diag}_m\cdot(W_0+\Mr)=\{\diag(\ell)(W_0+X):\rk X\le r\}$, exactly LoRA followed by a free per-output gain.
  Its dimension is $\min(mn,\,r(m+n-r)+m)$ for generic $W_0$ [Num], and smaller for special $W_0$, for instance
  $r(m+n-r)+m-r$ when $W_0$ has rank one. Its updates can exceed rank $r$. From $B_0=0$ the base point is an apex, as
  for LoRA.
\item[Fibre] The gauge contains $\GL_r$, acting on $(B,A)$ as for LoRA. For generic $W_0$ with $(m-r)(n-r)\ge m$ the
  budget $r(m+n)+m$ exceeds the dimension by exactly $r^2$, so $\GL_r$ is the whole continuous gauge. Gradient flow
  conserves the LoRA charge $\Phi$, also through the detached norm. The two backward rules give different geometries.
  The exact gradient of the row normalisation projects the direction gradient row by row, which is a change of metric.
  The detached backward of the DoRA paper (its Sec.~4.3) and of HF PEFT (\texttt{weight\_norm.detach()}) removes that
  projection and gives a non-gradient update, which need not be a descent direction. Under AdamW, which both the paper
  and HF use, neither variant follows a fixed metric, and the magnitude coordinates matter in their own right. Since $\mu_i$
  has size about $\lVert w_{0,i}\rVert$ and an output gain $\ell_i$ has size $1$, an equal step on each changes the
  row gain by different amounts.
\item[Merge] M1. Since the images agree, any gap between DoRA and LoRA plus an output gain, started at the same
  weights, comes from training dynamics. Experiment~D (\thmref{Proposition}{II.8}, \thmref{Prediction}{X.2}) separates
  its three possible sources, the parametrisation, the backward rule and the optimizer coordinates.
\end{exVIIlines}
\exVIIrefs{\thmref{Proposition}{II.8} \textperiodcentered{} \thmref{Proposition}{II.12} \textperiodcentered{}
  \thmref{Theorem}{III.5} \textperiodcentered{} \thmref{Prediction}{X.2}}
\end{example*}

% ------------------------------------------------------------------------------------------------ OFT / BOFT
\begin{example*}{OFT\,/\,BOFT}\phantomsection\label{ex:oft}%
\exVIIsrc{\citet{qiu2023controlling}; \citet{liu2023parameter}}{multiplicative}
\[
  W=W_0R\T\ \ \text{(OFT)},\qquad W=\diag(s)\,W_0R\T\ \ \text{(BOFT)},\qquad R_0=I,\ s_0=\mathbf 1.
\]
Here $R$ is block-diagonal with Cayley blocks $C(Q)=(I+Q)(I-Q)^{-1}$, $Q$ skew-symmetric of size $b$ (OFT), or a
product of $m'$ butterfly factors (BOFT).
\begin{exVIIlines}
\item[Image] With exact Cayley maps the image is an orbit in the level set of the neuron Gram matrix
  $K_{\rm out}=WW\T$, so every neuron norm, every angle between neurons and every singular value of $W_0$ stays fixed,
  and block OFT also fixes the $n/b$ partial Grams $W_{:,J}W_{:,J}\T$. BOFT has
  $K_{\rm out}(W)=\diag(s)K_{\rm out}(W_0)\diag(s)$. It keeps $\lvert\cos\rvert$ between neurons when every
  $s_i\ne0$, and the cosines themselves when all $s_i$ share a sign; with one butterfly factor it sends each block Gram
  to its $\diag(s)$-rescaling. HF's default Cayley--Neumann OFT, $R=C(Q)(I-Q^4)$, is only approximately orthogonal,
  since $R\T R=(I-Q^4)^2$. The base point is regular. For BOFT this holds when the factors' plane sets are disjoint
  (for instance $b=2$), so that the factor Lie algebras form a direct sum.
\item[Fibre] For exact-Cayley block OFT the gauge is the stabiliser of $W_0$ in the block group, which is trivial
  when $W_0$ is injective.
\item[Base] Exact-Cayley block OFT on a fixed partition is idempotent up to closure, and stays in
  $W_0\cdot\prod\SO(b)$ however often it is merged. Iterated BOFT with $m'$ factors reaches
  $\diag\cdot W_0\cdot P\big(\prod\SO(b\,2^{m'-1})\big)P\T$ for a fixed permutation $P$, and, for injective $W_0$,
  $\diag\cdot W_0\cdot\SO(n)$ only at full butterfly depth. At $n=64$
  and $b=4$ the block hypergraph has components of sizes $4,8,16,32,64$ for $m'=1,\dots,5$, so it is connected only at
  $m'=5$ [Num] (Figure~\ref{fig:rebase}b). With small generators ($\lVert Q\rVert_{\rm op}\le2^{1/4}$) each merge of
  HF's default is a contraction, so re-merging can only shrink neuron norms and singular values.
\item[Merge] M1.
\end{exVIIlines}
\exVIIrefs{\thmref{Theorem}{II.5} \textperiodcentered{} \thmref{Corollary}{II.6} \textperiodcentered{}
  \thmref{Theorem}{V.3}}
\end{example*}

% ------------------------------------------------------------------------------------------------ (IA)^3
\begin{example*}{(IA)$^3$}\phantomsection\label{ex:ia3}%
\exVIIsrc{\citet{liu2022few}}{multiplicative}
\[
  W=\diag(\ell)\,W_0\ \ (W_K,\,W_V),\qquad W=W_0\,\diag(\ell)\ \ (W_{\rm down}),\qquad \ell_0=\mathbf 1.
\]
\begin{exVIIlines}
\item[Image] Linear, the closure of the orbit of $W_0$ under a torus of one-sided diagonal scalings. Zeros of $W_0$
  stay zero, and row lines (column lines on $W_{\rm down}$) stay fixed while all $\ell_i\ne0$. The image lies inside
  that of rank-1 HiRA, $\mathrm{HiRA}_1$. Its covariance group is $\mathrm{Mon}\times\GL$ on the output-scaled boxes
  and $\GL\times\mathrm{Mon}$ on $W_{\rm down}$.
\item[Base] Idempotent.
\item[Merge] M1$\to$. It fuses into $W_K$, $W_V$ and $W_{\rm down}$ for every activation. The FFN scale fuses
  alternatively into $W_{\rm up}$ under SwiGLU or GeGLU for every $\ell$, through the linear branch, a scale migration
  that AWQ already uses; under ReLU for $\ell\ge0$; and under GELU only for entries $\ell_i\in\{0,1\}$.
\end{exVIIlines}
\exVIIrefs{\thmref{Theorem}{II.5} \textperiodcentered{} \thmref{Proposition}{II.9} \textperiodcentered{}
  \thmref{Theorem}{VI.3}}
\end{example*}

% ------------------------------------------------------------------------------------------------ VeRA
\begin{example*}{VeRA}\phantomsection\label{ex:vera}%
\exVIIsrc{\citet{kopiczko2023vera}}{additive}
\[
  \rho(b,d)=W_0+\Lambda_bB\Lambda_dA,\qquad q_0=(0,d_0).
\]
Here $\Lambda_b=\diag(b)\in\R^{m\times m}$ and $\Lambda_d=\diag(d)\in\R^{r\times r}$ are trained, while
$B\in\R^{m\times r}$ and $A\in\R^{r\times n}$ are frozen, random and shared across layers.
\begin{exVIIlines}
\item[Image] A toric cone of dimension $m+r-1$ [Num], with apex at $\theta_0$. At $b_0=0$ only $b$ receives a
  gradient, so the first-order image has dimension at most $m$, below $m+r-1$ when $r\ge2$.
\item[Fibre] The gauge is $\R^\times$, with $c$ acting by $(b,d)\mapsto(cb,d/c)$, and gradient flow conserves the
  charge $\lVert b\rVert^2/\eta_b-\lVert d\rVert^2/\eta_d$.
\item[Base] If $B$ has no zero entry, rebasing reaches exactly LoRA-FA$(A)$, the method $W_0+\{XA\}$, which is the
  target of the arrow $h(b,d)=\Lambda_bB\Lambda_d$ from VeRA to LoRA-FA$(A)$.
\item[Merge] M1.
\end{exVIIlines}
\exVIIrefs{\thmref{Lemma}{II.4} \textperiodcentered{} \thmref{Proposition}{II.12} \textperiodcentered{}
  \thmref{Theorem}{III.5} \textperiodcentered{} \thmref{Theorem}{V.3}}
\end{example*}

% ------------------------------------------------------------------------------------------------ LoHa / LoKr
\begin{example*}{LoHa\,/\,LoKr}\phantomsection\label{ex:loha}%
\exVIIsrc{\citet{hyeonwoo2021fedpara}; \citet{yeh2023navigating}}{additive}
\[
  W_0+(B_1A_1)\odot(B_2A_2)\ \ \text{(LoHa)},\qquad W_0+C\otimes BA\ \ \text{(LoKr)}.
\]
\begin{exVIIlines}
\item[Image] LoHa$_{r,r}$ has updates of rank at most $\min(r^2,m,n)$, by the cut-rank bound. Its dimension is at most
  $\min\big(mn,\,2r(m+n)-2r^2-(m+n-1)\big)$, with equality below $mn$ [Num]. The gauge is $\GL_r\times\GL_r$ times the
  two-sided diagonal torus $(B_1,A_1,B_2,A_2)\mapsto(D_1B_1,A_1D_2,D_1^{-1}B_2,A_2D_2^{-1})$, and their joint generic
  orbits have dimension $2r^2+m+n-1$. HF zero-initialises one factor, so the base point is an apex at which only that
  factor moves to first order. At the common budget $2r(m+n)$ of LoHa$_{r,r}$ and LoRA$_{2r}$ three regimes occur. For
  $r\le2$, $r^2\le2r$, so the image of LoHa$_{r,r}$ lies inside that of LoRA$_{2r}$, strictly when $2r^2<m+n-1$. For
  $r\ge3$ and $2r^2<m+n-1$, LoHa trades at least $(m+n-1)-2r^2$ dimensions for rank up to $r^2$. When
  $2r^2\ge m+n-1$, for instance at $r=64$ on $4096\times4096$ or at LyCORIS-style $r=32$ on $320$- to $640$-wide
  layers, it matches or exceeds LoRA$_{2r}$ in dimension as well as in rank. At $131{,}072$ parameters on a
  $4096\times4096$ matrix, LoRA$_{16}$ reaches $130{,}816$ dimensions and LoHa$_{8,8}$ at most $122{,}753$
  (Figure~\ref{fig:rankvol}).
\item[Image] A $k$-term Kronecker sum $\sum_iC_i\otimes D_i$ is a matrix of rank at most $k$ after the Van
  Loan--Pitsianis rearrangement $\mathcal R$, which sends $C\otimes D$ to $\vect(C)\vect(D)\T$. So LoKr is a
  constrained LoRA$_1$ transported by $\mathcal R^{-1}$, which is LoRA across a different cut of the same four-leg
  tensor. Write $m=m_1m_2$ and $n=n_1n_2$. For $C\in\R^{m_1\times n_1}$ and inner rank $r'\le\min(m_2,n_2)$, LoKr's
  maximum rank is $\min(m_1,n_1)\,r'$, attained. It lies inside zero-init LoRA$_r$ iff $r\ge\min(m_1,n_1)\,r'$, and
  is otherwise incomparable with it when $\min(m_1,m_2)\min(n_1,n_2)>1$ (Figure~\ref{fig:cut}). In HF PEFT the second
  factor is low-rank only when $r<\max(m_2,n_2)/2$; above that LoKr is an unconstrained Kronecker product.
\item[Merge] M1.
\end{exVIIlines}
\exVIIrefs{\thmref{Theorem}{II.10} \textperiodcentered{} \thmref{Theorem}{II.11} \textperiodcentered{}
  \thmref{Proposition}{II.12} \textperiodcentered{} \thmref{Proposition}{II.13}}
\end{example*}

% ------------------------------------------------------------------------------------------------ FourierFT
\begin{example*}{FourierFT}\phantomsection\label{ex:fourierft}%
\exVIIsrc{\citet{gao2024parameter}}{additive}
\[
  W_0+\alpha\,\operatorname{Re}F^{-1}(c),\qquad c\ \text{supported on $n_f$ fixed frequencies}.
\]
Here $F$ is the two-dimensional discrete Fourier transform and $\alpha$ a fixed scaling (HF's default $150$).
\begin{exVIIlines}
\item[Image] A coordinate mask transported into the real cosine frame. Its image lies in the point-symmetric
  subspace $\Delta W_{j,l}=\Delta W_{-j,-l}$. Started at $\theta_0$ ($c_0=0$, HF's \texttt{init\_weights=True}) it is linear
  and regular. Its dimension is $n_f$ minus the number of sampled conjugate pairs $(u,v),(-u,-v)$, each of which adds
  one gauge dimension. At LLM sizes this is close to $n_f$, but at $768\times768$ with $1000$ frequencies a pair occurs
  in about half the draws, and HF's default seed $777$ gives exactly two. HF's default random spectrum
  ($\mathcal N(0,1)$) starts elsewhere, with a defect about $14\%$ the size of a $768\times768$ $W_0$ of entry standard
  deviation $0.04$ at $1000$ frequencies. Its covariance group $\Pi_{\rm cyc}$ excludes most neuron permutations, so
  permuting hidden units can change what it reaches.
\item[Base] Idempotent with fixed indices.
\item[Merge] M1.
\end{exVIIlines}
\exVIIrefs{\thmref{Proposition}{II.12} \textperiodcentered{} \thmref{Proposition}{II.13} \textperiodcentered{}
  \thmref{Observation}{II.15} \textperiodcentered{} \thmref{Theorem}{V.3}}
\end{example*}

% ------------------------------------------------------------------------------------------------ BitFit
\begin{example*}{BitFit}\phantomsection\label{ex:bitfit}%
\exVIIsrc{\citet{benzaken2021bitfit}}{selective}
\[
  b\mapsto b_0+\beta\ \ \text{on every bias},\qquad W=W_0.
\]
\begin{exVIIlines}
\item[Image] A translation of the bias block, linear and regular. It has no gauge, so its expressive dimension equals
  its parameter count.
\item[Base] Idempotent.
\item[Merge] Nothing to merge, since BitFit trains the model's own biases. Llama and Mistral have no biases, so it
  trains nothing there.
\end{exVIIlines}
\exVIIrefs{\thmref{Proposition}{II.12} \textperiodcentered{} \thmref{Theorem}{V.3}}
\end{example*}

% ------------------------------------------------------------------------------------------------ Adapters
\begin{example*}{Adapters}\phantomsection\label{ex:adapter}%
\exVIIsrc{\citet{houlsby2019parameter}}{architectural}
\[
  x\mapsto x+U\sigma(Dx),\qquad D\in\R^{r\times d},\ U\in\R^{d\times r}.
\]
\begin{exVIIlines}
\item[Image] Function level. The adapter is an extension of the network, with neutral point $U=0$ (zero-up
  initialisation). Houlsby et al.'s own near-identity initialisation, with weights drawn from a zero-mean Gaussian of
  standard deviation $10^{-2}$ truncated at two standard deviations, starts near the neutral point and leaves a small
  defect.
\item[Fibre] Under Houlsby's GELU the gauge is only the permutations of hidden units. A ReLU, as in the Pfeiffer and
  parallel adapters, adds positive per-unit rescalings, giving $S_r\ltimes\R_{>0}^r$. At $U=0$ only $U$ moves to first
  order (the zero gate, function-level variant).
\item[Merge] M$\infty$, since it is not box-locally mergeable for non-linear $\sigma$. With $\sigma=\id$, an adapter parallel to a
  single linear box is LoRA, with the same $\rho$; one parallel to a whole sublayer edits the weightless skip
  connection and stays M$\infty$ when the sublayer is non-affine. A serial linear adapter $(I+UD)W_0$ after a
  bias-free $W_0$ reaches exactly the rank-$\le r$ updates whose rows lie in $\operatorname{row}W_0$, which is LoRA's
  image iff $\rk W_0=n$.
\end{exVIIlines}
\exVIIrefs{\thmref{Lemma}{II.4} \textperiodcentered{} \thmref{Definition}{VI.1} \textperiodcentered{}
  \thmref{Lemma}{VI.2} \textperiodcentered{} \thmref{Theorem}{VI.3} \textperiodcentered{}
  \thmref{Proposition}{VI.4}}
\end{example*}

% ------------------------------------------------------------------------------------------------ Prefix / prompt
\begin{example*}{Prefix\,/\,prompt tuning}\phantomsection\label{ex:prefix}%
\exVIIsrc{\citet{li2021prefix}; \citet{lester2021power}}{augmentation}
\[
  \mathrm{Attn}_q(P\uplus C)=(1-\lambda_q)\,\mathrm{Attn}_q(C)+\lambda_q\,\mathrm{Attn}_q(P),\qquad
  \lambda_q=\frac{Z_q(P)}{Z_q(P)+Z_q(C)}.
\]
Here $P$ and $C$ are the prefix and content multisets of key--value pairs, and
$Z_q(K)=\sum_{k\in K}e^{\langle q,k\rangle/\sqrt d}$.
\begin{exVIIlines}
\item[Image] Prefix and prompt tuning are unpointed extensions on the sequence wire, in the sense of
  \thmref{Theorem}{VI.5}(f): on a head with $W_oW_v\ne0$ no prefix is neutral. A prefix acts as an input-gated
  parallel adapter on attention. Within its layer it multiplies the content attention weights by $1-\lambda_q$ and
  keeps their ratios, and the output moves on the segment towards $\mathrm{Attn}_q(P)\in\operatorname{conv}V_P$
  (Figure~\ref{fig:prefix}). This rigidity is relative to the same layer without the prefix, and relative to the
  pretrained pattern only at the lowest prefixed layer.
\item[Arrows] In causal decoders whose content positions are offset by the prompt length, as in HF PEFT, every soft
  prompt is realised exactly by the deep prefix made of its own per-layer keys and values. This is an arrow of the
  unpointed slice, from prompt tuning into deep prefix tuning (P-tuning v2). \citet{petrov2023prompting} and
  \citet{li2021prefix} assert that the inclusion is strict. It is strict in a one-layer, one-token example, where a
  prompt fixes the prefix value as a function of its key; the generic statement is not proved here, and encoders remain
  open. LLaMA-Adapter \citep{zhang2023llama} normalises the prefix separately and adds a zero-initialised gate, so it
  leaves the content attention weights unchanged and supplies the missing neutral point.
\item[Merge] M$\infty$. For a standard softmax head with $W_o\ne0$ and $W_q\ne0$, a generic prefix makes the head
  non-affine even on single-token inputs, so prefix tuning cannot be folded into the head's own weights.
\end{exVIIlines}
\exVIIrefs{\thmref{Definition}{VI.1} \textperiodcentered{} \thmref{Theorem}{VI.5}}
\end{example*}

% ------------------------------------------------------------------------------------------------ PiSSA / LoftQ / QLoRA
\begin{example*}{PiSSA\,/\,LoftQ\,/\,QLoRA}\phantomsection\label{ex:pissa}%
\exVIIsrc{\citet{meng2024pissa}; \citet{li2023loftq};\newline \citet{dettmers2023qlora}}{additive}
\begin{gather*}
  \rho(B,A)=W_{\rm res}+sBA,\qquad sB_0A_0=[W_0]_r,\qquad W_{\rm res}=W_0-[W_0]_r\qquad\text{(PiSSA)},\\
  \rho(B,A)=\kappa\theta_0+sBA,\qquad B_0=0\qquad\text{(QLoRA)}.
\end{gather*}
Here $[X]_r$ is the best rank-$r$ approximation of $X$, $W_{\rm res}$ is frozen, and $\kappa$ quantises and
dequantises. \citet{meng2024pissa} take $s=1$.
\begin{exVIIlines}
\item[Image] PiSSA is a split (Type II) start. Its image is the cone with vertex $W_0-[W_0]_r$, and $\theta_0$ is a
  smooth point of it (Figure~\ref{fig:apex}b). Its image and zero-init LoRA's are incomparable. No exact pointing of
  LoRA$_r$, zero-init or split, makes an update of rank above $2r$ reachable, and split pointings attain $2r$ when
  $2r\le\min(m,n)$.
\item[Fibre] Gauge $\GL_r$. HF's PiSSA factors give $\Phi_0=(S_r/s)(1/\eta_B-1/\eta_A)$, so the start is balanced,
  $\Phi_0=0$, at equal rates. By the balance--escape trilemma a balanced, non-stationary, identity-at-init start needs
  the frozen residual.
\item[Base] QLoRA is LoRA pointed at $\kappa\theta_0$, with generically full-rank defect $e=\kappa\theta_0-\theta_0$,
  so its image misses $\theta_0$ unless $\rk e\le r$ (Figure~\ref{fig:teaser}). Splitting and error fitting change
  the defect, and the order of operations matters. QPiSSA splits, then quantises the residual, and leaves
  $\kappa(W_{\rm res})-W_{\rm res}$. One-step LoftQ quantises, then fits the error. HF's LoftQ initialisation ignores
  the LoRA scale $s=\alpha/r$, so the defect is $e-s[e]_r$. It is $e$ minus its best rank-$r$ approximation only when
  $s=1$, and for $s=2$, as with $\alpha=2r$, its norm equals $\lVert e\rVert$.
\item[Merge] PiSSA merges into its box (M1). On a quantised base the type is Mq. Exact merges of generic
  high-precision adapters leave the quantised type, and type-preserving merges are lossy; QA-LoRA's group-constant
  adapters under affine group quantisation are the exception \citep{xu2023qa}.
\end{exVIIlines}
\exVIIrefs{\thmref{Theorem}{II.1} \textperiodcentered{} \thmref{Theorem}{II.2} \textperiodcentered{}
  \thmref{Corollary}{III.6} \textperiodcentered{} \thmref{Proposition}{III.8} \textperiodcentered{}
  \thmref{Proposition}{V.2}}
\end{example*}

% ------------------------------------------------------------------------------------------------ ReLoRA
\begin{example*}{ReLoRA}\phantomsection\label{ex:relora}%
\exVIIsrc{\citet{lialin2023relora}}{optimizer}
\[
  \theta_{k+1}=\theta_k+sB_kA_k,\qquad \text{fresh }(B_k,A_k)\text{ at each stage}.
\]
\begin{exVIIlines}
\item[Base] Rebased LoRA. After $K$ stages it reaches exactly $W_0+\MM_{\le\min(Kr,m,n)}$, so all of $\Theta$ at
  $K^*=\lceil\min(m,n)/r\rceil$ (Figure~\ref{fig:rebase}a). Over weight space it is $\mathrm{LoRA}_r^{\oplus K}$
  unrolled in time.
\item[Fibre] Each stage is an apex with gauge $\GL_r$.
\item[Merge] Each stage merges (M1).
\end{exVIIlines}
\exVIIrefs{\thmref{Observation}{I.7} \textperiodcentered{} \thmref{Theorem}{V.3}}
\end{example*}

% ------------------------------------------------------------------------------------------------ GaLore
\begin{example*}{GaLore}\phantomsection\label{ex:galore}%
\exVIIsrc{\citet{zhao2024galore}}{optimizer}
\[
  W_{t+1}=W_t-\eta\,P_t\,\mathrm{opt}\big(P_t\T G_t\big),\qquad P_t\in\R^{m\times r}.
\]
Here $G_t$ is the weight gradient, and $P_t$ is refreshed from its SVD every $T$ steps.
\begin{exVIIlines}
\item[Fibre] A backward lens with compression $P_t\T$ and anchor $P_t$. Under an optimizer that reads only its
  gradients, such as Adam with weight decay $0$ and no full-gradient clipping, each period is exactly one-sided LoRA
  $W_k+P_kC$ with $C_0=0$, merged and re-initialised at the end of the period, with its optimizer state carried over
  rather than reset or transported. Decoupled weight decay and full-gradient clipping break this identity.
\item[Base] The anchor is constant on each period, so it is integrable trivially, including at $T=1$, and
  integrability does not separate GaLore from LoRA. What separates them is the resampling schedule together with the
  carried state. As a set over its possible sequences of projectors, the moving shape reaches all of $\Theta$.
\item[Merge] Not applicable, since GaLore writes $W$ itself.
\end{exVIIlines}
\exVIIrefs{\thmref{Definition}{IV.1} \textperiodcentered{} \thmref{Proposition}{IV.2} \textperiodcentered{}
  \thmref{Proposition}{IV.3} \textperiodcentered{} \thmref{Theorem}{V.3}}
\end{example*}

% ------------------------------------------------------------------------------------------------ Merging
\begin{example*}{Merging}\phantomsection\label{ex:merging}%
\exVIIsrc{\citet{ilharco2022editing}; \citet{huang2023lorahub}}{composition}
\[
  \theta_0+\textstyle\sum_i\lambda_i\tau_i\ \ \text{(task arithmetic)},\qquad
  \big(\textstyle\sum_iw_iB_i\big)\big(\textstyle\sum_jw_jA_j\big)\ \ \text{(LoraHub)}.
\]
\begin{exVIIlines}
\item[Fibre] Task arithmetic is $\GL(\Theta)$-natural. TIES-Merging \citep{yadav2023ties}, with global trimming and
  sign election by total mass, is natural under homotheties times signed permutations. DARE \citep{yu2023language} is
  natural under every invertible coordinate rescaling for each mask and under permutations in law, so its
  distributional group is the monomial group. None of these sparsifying rules is natural under rotations.
\item[Fibre] Post-hoc factor-wise rules (LoraHub, and HF's \texttt{linear} merge with its non-SVD \texttt{ties} and
  \texttt{dare} variants) depend on each adapter's gauge, since replacing $(B_2,A_2)$ by $(cB_2,A_2/c)$ changes
  LoraHub's result by $\sum_{l\ne2}w_lw_2\big[(c-1)B_2A_l+(c^{-1}-1)B_lA_2\big]$ (Figure~\ref{fig:gauge}). A shared
  factor, such as a common frozen $A$, or a joint frame that produces one, such as the joint $U\Sigma$ of KnOTS
  \citep{stoica2024model}, removes the gauge dependence. With a shared $A$, LoraHub returns
  $(\sum_jw_j)\sum_iw_iB_iA$, exact up to the global factor $\sum_jw_j$.
\item[Merge] Merging an adapter into a full-precision base is the canonical arrow to the terminal object, full
  fine-tuning.
\end{exVIIlines}
\exVIIrefs{\thmref{Observation}{I.4} \textperiodcentered{} \thmref{Proposition}{V.4}}
\end{example*}

Exposé~\ref{sec:atlas} turns the answers on these cards into seven coordinates, each read off a method's formula
without training, and into the arrows between methods.
